\begin{table}[!ht]
\centering
\caption{Decomposição RIF-OB (Firpo, Fortin \& Lemieux, 2018) do gap salarial racial
por quantil incondicional, em formato \emph{two-fold} (referência: estrutura de preços
dos brancos). Dotações: diferença de características observáveis (capital humano, ocupação,
contexto). Retornos: parcela não explicada (componente discriminatório). Dotações + Retornos
$=100\%$ do gap em cada quantil. Padrão \emph{sticky floor}: o componente de retornos
(discriminação de mercado) é maior na base e decresce rumo ao topo. PNAD Contínua 2016--2025.}
\label{tab:rif_ob}
\begin{tabular}{lrrr}
\toprule
Quantil & Gap obs. & Dotações (\%) & Retornos (\%) \\
\midrule
q10 & 0,559 & 64,9 & 35,1 \\
q25 & 0,310 & 68,8 & 31,2 \\
q50 & 0,378 & 72,7 & 27,3 \\
q75 & 0,470 & 75,5 & 24,5 \\
q90 & 0,483 & 87,1 & 12,9 \\
\bottomrule
\end{tabular}
\end{table}
